% !PRINTED 281
\weberpage{281}
%<<CONT
zu untersuchenden Ueberschiebungen mit $f$, $(f, f)^2$, $((f, f)^2, f)$ übrig.
Nun ist bereits $(ak)^3$ Reducent, und zwar ist $(ak)^3 Q \equiv 0 \bmod (ak)^4$,
wo $Q$ irgend ein symbolisches Product. Folglich sind höhere als
zweite Ueberschiebungen überflüssig, und wir können uns demnach
auf die erste Potenz von $k$ beschränken.

Durch ein- und zweimalige Ueberschiebung von $k$ über die
Formen des Systemes $(A^{(1)})$ entstehen aber die Covarianten:

\[
\begin{aligned}
&(f,\, k), \quad (f,\, k)^2, \quad ((f,\, f)^2,\, k); \quad ((f,\, f)^2,\, k)^2, \quad \bigl(((f,\, f)^2,\, f),\, k\bigr),\\
&\qquad \bigl(((f,\, f)^2,\, f),\, k\bigr)^2.
\end{aligned}
\]

Die drei letzten Formen sind überflüssig; denn $((f, f)^2, k)^2$ ist schon
nach Nr. 255 (II) als eine reducible Form erkannt worden. Die Form
$\bigl(((f, f)^2, f), k\bigr)$ ist aber reducibel als Functionaldeterminante von einer
Functionaldeterminante und im gleichen Sinne ist auch $\psi = \bigl(((f,f)^2, f), k\bigr)^2$
eine reducible Form, in so ferne die Glieder dieser Ueberschiebung
durch den Productsatz in Aggregate umgestaltet werden können, deren
einzelne Terme theils reducible, theils bereits in das System aufgenommene
 Formen zu Factoren haben. Hievon überzeugt man sich
unmittelbar auch auf folgende Weise. Gemäss den Ueberschiebungsgesetzen
 Nr. 44 ist, da $(f, f)^5 = k^{(1)} = 0$,

\[
\bigl(((f, f)^2, f), k\bigr)^2 = G_0 + \sum{}'(T^{(1)}, k)^1 + \sum{}'(T^{(2)}, k)^0 + \sum{}'(T, k^{(2)})^0,
\]

wobei das Anfangsglied $G_0 = (ab)^2 (ak)^2 (ac) a_x b_x^4 c_x^5 k_x^2$ und $T^{(1)}, T^{(2)}$;
$k^{(1)}, k^{(2)}$ Formen sind, die aus $T = ((f, f)^2, f)$ resp. $k = (f, f)^4$ durch
einmalige, resp. zweimalige Faltung entstehen. Die letzte Summe
rechts hat nur den einen Summanden $T \cdot A$; alle Summanden der
zweiten Summe zerfallen in Producte, deren einer Factor die Form $k$.
Die Formen $T^{(1)}$ in der ersten Summe sind

\[
\begin{aligned}
(ab)^3 (ac) a_x^2 b_x^3 c_x^5 &= \tfrac{1}{2} (ab)^3 a_x^2 b_x^2 c_x^5 \{(ac) b_x - (bc) a_x\} = \tfrac{1}{2} k \cdot f\\[4pt]
(ab)^2 (ac)^2 a_x^2 b_x^4 c_x^4 &= \tfrac{1}{4} \{(ab)^2 c_x^2 + (ac)^2 b_x^2 - (bc)^2 a_x^2\}^2 a_x^2 b_x^2 c_x^2 = \tfrac{1}{2} k \cdot f\\[4pt]
(ab)^2 (ac) (bc) a_x^2 b_x^3 c_x^4 &= \tfrac{1}{2} \{(ac)^2 b_x^2 + (bc)^2 a_x^2 - (ab)^2 c_x^2\} (ab)^2 a_x^2 b_x^2 c_x^2 = 0,
\end{aligned}
\]

so dass also auch die Formen $(T^{(1)}, k)^1$, wenn sie nicht verschwinden,
eine der Formen $k$ oder $f$ zu Factoren haben. $G_0$ ist aber ohnehin
reducibel nach Nr. 255 und damit die ganze rechte Seite.


% ===== pdf idx 296, printed 282.0 =====
% !PRINTED 282
\weberpage{282}
Das ganze System $S \equiv 0 \bmod (ab)^4$ besteht demnach aus den
zehn Formen:

\[
\left.
\begin{array}{l}
f, \quad (f, f)^2, \quad ((f, f)^2, f)\\[2pt]
k, \quad (f, k), \quad (f, k)^2, \quad ((f, f)^2, k)\\[2pt]
A, \quad B, \quad C.
\end{array}
\right\}\ S
\]

260. \textit{Ueberschiebung dieses Systemes $S$ über $l = (ak)^4 a_x^2$.} Unsere
letzte Aufgabe ist somit, die Ueberschiebungen der sieben Covarianten
dieses Systemes über die Covariante

\[
l = (ak)^4 a_x^2
\]

zu studiren. Dabei ist es überhaupt nicht nöthig, Producte dieser
Formen über $l$ zu schieben. Denn alle sieben Covarianten sind geraden
 Grades, und demnach würden in jeder Ueberschiebung eines Productes
 dieser Formen über Potenzen von $l$ zerfallende Glieder gebildet
 werden können.

Wir haben also nur die erste Potenz jeder der sieben Covarianten
über Potenzen von $l$ zu schieben.
Die Form $f$ liefert dadurch die sechs neuen Formen

\[
(f, l), \quad (f, l)^2, \quad (f, l^2)^3, \quad (f, l^2)^4, \quad (f, l^3)^5, \quad (f, l^3)^6.
\]

Die Form $(f, f)^2$ giebt nur zu einer neuen Form

\[
((f, f)^2, l)
\]

Veranlassung; denn bereits bei der zweiten Ueberschiebung tritt der
Reducent $(ab)^2 (al)^2$ auf. [Vgl. Nr. 257 (IV)].

Aus den Ueberschiebungen von $((f, f)^2, f)$ mit $l$ entstehen überhaupt
 keine neuen Formen; denn die erste ist als Functionaldeterminante
 von einer Functionaldeterminante reducibel; alle weiteren besitzen
 den Reducenten $(al)^2 (ab)^2$.

Als weitere Formen des Systemes ergeben sich ferner die Ueberschiebungen


\[
(k, l), \quad (k, l)^2, \quad (k, l^2)^3,
\]

während die Invariante $(k, l^2)^4$ nach (V) Nr. 258 reducibel ist.

Von den Ueberschiebungen der Form $(f, k)$ über Potenzen von $l$
liefern die ungeraden überflüssige Formen. Denn sie enthalten resp.
Glieder wie:

\[
\begin{aligned}
&(ak) (al) l_x a_x^4 k_x^3\\
&(ak) (al)^2 (al_1) l_x a_x^2 k_x^3\\
&(ak) (al)^2 (al_1)^2 (al_2) l_{2x} k_x^3\\
&(ak) (al)^2 (al_1)^2 (al_2) (kl_3)^2 l_{2x} k_x,
\end{aligned}
\]
%CONT>>


% ===== pdf idx 297, printed 283.0 =====
% !PRINTED 283
\weberpage{283}
%<<CONT
welche stets wegen des Productsatzes auf andere Formen des Systemes
reducibel sind. So ist z. B.

\[
\begin{aligned}
(ak) (al) k_x l_x \cdot a_x^4 k_x^2 &= a_x^4 k_x^2 \{(ak)^2 l_x^2 + (al)^2 k_x^2 - (kl)^2 a_x^2\}\\
&= (f, k)^2 \cdot l + (f, l)^2 \cdot k - (f, l)^2 \cdot (k, l)^2.
\end{aligned}
\]

Dagegen liefern die geraden Ueberschiebungen vier neue Formen:

\[
((f, k), l)^2, \quad ((f, k), l^2)^4, \quad ((f, k), l^3)^6, \quad ((f, k), l^4)^8.
\]

Hiezu tritt als weitere Covariante die erste Ueberschiebung der sechsten
Covariante $(f, k)^2$ über $l$, also die Form

\[
((f, k)^2, l),
\]

während bereits ihre zweite Ueberschiebung wegen des Reducenten
$(ak)^2 (al)^2$ reducibel ist.

Die siebente Covariante $((f, f)^2, k)$ endlich liefert keine neue Form
mehr. Denn ihre erste Ueberschiebung mit $l$ ist nach dem nun wiederholt
 citirten Satze als Functionaldeterminante reducibel, und jede
weitere hat den Reducenten $(ab)^2 (al)^2$.

261. \textit{Tabelle der Formen des Systemes $(A^{(2)})$.} Die Gesammtheit
aller Formen des vollen Systemes von $f = a_x^6$ ist also durch folgende
26 Formen gegeben.

1) vier gerade Invarianten:

\[
A, \ B, \ C, \ (f, l^3)^6 \text{ und eine schiefe } ((f, k), l^4)^8;
\]

2) sechs quadratische Covarianten:

\[
l, \ (k, l)^2, \ (f, l^2)^4, \ (k, l^2)^3 \ (f, l^3)^5, \ ((f, k), l^3)^6;
\]

3) fünf biquadratische Formen:

\[
k, \ (f, l)^2, \ (k, l), \ (f, l^2)^3, \ ((f, k), l^2)^4;
\]

4) fünf Covarianten sechsten Grades in $x$:

\[
f, \ (f, k)^2, {}^{*)} \ (f, l), \ ((f, k)^2, l), \ ((f, k), l)^2;
\]

5) drei Covarianten achten Grades:

\[
(f, f)^2, \ (f, k), \ ((f, f)^2, l);
\]

6) eine Covariante zehnten und eine zwölften Grades, nämlich:

\[
((f, f)^2, k), \ \text{resp.} \ ((f, f)^2, f).
\]

\bigskip
\noindent\textbf{§ 27. Relationen zwischen den Formen des Systemes von $f = a_x^6$.}
\bigskip

262. \textit{Einführung von Ueberschiebungsgliedern in das System.} Nachdem
 wir durch die vorausgegangenen Untersuchungen in den Besitz
sämmtlicher Formen des vollen Systemes von $f = a_x^6$ gelangt sind,
%CONT>>

\bigskip
\noindent\rule{4cm}{0.4pt}

\noindent{*)} $(f, k)^2 = p$ bei \gesperrt{Clebsch} \glqq Binäre Formen\grqq.


% ===== pdf idx 298, printed 284.0 =====
% !PRINTED 284
\weberpage{284}
%<<CONT
tritt an uns die Aufgabe heran, diese Formen, wenn nöthig, durch
eingliedrige symbolische Producte zu ersetzen, die für die symbolische
Rechnung als besonders tauglich erscheinen. Dies gilt insbesondere
für die fünf Invarianten des Systemes, sowie für die sechs einfachsten
Covarianten, d. i. für die sechs quadratischen Formen.

Von den Invarianten sind aber die drei ersten ohnehin einfache
symbolische Producte, nämlich

\[
\begin{aligned}
A &= (ab)^6, \quad B = (kk_1)^4,\\
C &= (kk_1)^2 (kk_2)^2 (k_1k_2)^2;
\end{aligned}
\]

diese werden wir beibehalten. Dagegen werden wir an Stelle der
Ueberschiebungen $(f, l^3)^6$, $((f, k), l^4)^8$ symbolische Producte einführen,
die sich aus dem simultanen System der sechs quadratischen Covarianten
 ergeben.

Von den quadratischen Covarianten sind die beiden Formen $l$
und $(k, l)^2$ durch je ein einfaches symbolisches Product dargestellt;
alle übrigen ersetzen wir durch eines ihrer Glieder.

263. \textit{Die quadratischen Formen.} Wie bei der Form fünften
Grades die linearen Covarianten, so spielen hier die quadratischen eine
wichtige Rolle. Auch hier ist das simultane System von vornherein
bekannt, und die erste Aufgabe ist, dessen Stellung zum vollen System
von $f = a_x^6$ kennen zu lernen. Diese Aufgabe aber deckt sich mit der
vorhin gestellten: die complicirteren quadratischen Covarianten durch
einfache symbolische Producte zu ersetzen. Setzen wir nämlich neben

\weq{1}{(ak)^4 a_x^2 = l, \quad (k, l)^2 = (kl)^2 k_x^2 = m,}

die wir a priori beibehalten, auch noch

\[
(k, m)^2 = n,
\]

dann zeigt sich, dass die in dem Systeme dieser drei Formen $l, m, n$
auftretenden Functionaldeterminanten

\weq{2}{\lambda = (m, n), \quad \mu = (n, l), \quad \nu = (l, m)}

Glieder von $(k, l^2)^3$ resp. $(f, l^3)^5$, $((f, k), l^3)^6$ sind, während $n$ an
Stelle der Ueberschiebung $(f, l^2)^4$ treten kann. Denn man hat wegen
der Relation [Nr. 258 (I)]

\[
(f, l)^2 = 2\varDelta + \tfrac{1}{3} A \cdot k
\]
% checked against the GDZ scan (27 Oct 2026): the usual slanted capital Delta of this fount, \varDelta

durch zweimalige Ueberschiebung mit $l$ die Beziehung

\[
(f, l^2)^4 = 2(\varDelta, l)^2 + \tfrac{1}{3} A (k, l)^2.
\]

Hierin ist nach unseren Bezeichnungen (1) $(k, l)^2 = m$; das erste
%CONT>>


% ===== pdf idx 299, printed 285.0 =====
% !PRINTED 285
\weberpage{285}

Glied rechts ergiebt sich aber als zweite Ueberschiebung von $\Delta=(k,\ k)^2$
über $l$ aus der Reihenentwicklung

\[
(kk_1)^2 k_x^2 k_{1y}^2 = \Delta_{y^2} + \frac{1}{3} B (yx)^2 ,
\]

indem wir darin $y$ durch $l$ ersetzen; man erhält

\[
(kk_1)^2 (k_1 l)^2 k_x^2 = (\Delta,\ l)^2 + \frac{1}{3} B \cdot l_x^2
\]

oder

\[
\big((k,\ l)^2,\ k\big)^2 = (\Delta,\ l)^2 + \frac{1}{3} B \cdot l
\]

oder

\weq{3}{(m,\ k)^2 = n = (\Delta,\ l)^2 + \frac{1}{3} B \cdot l\,.}
Demnach wird:

\weq{4}{(f,\ l^2)^4 = 2n - \frac{2B}{3}\cdot l + \frac{A}{3} m\,,}
und folglich kann $(f,\ l^2)^4$ durch $n$ ersetzt werden.

Es ist ferner:

\[
\begin{aligned}
(k,\ l^2)^3 &= \frac{2}{2} (kl)^2(kl_1)k_xl_x = \big((k,\ l)^2,\ l\big) = (m,\ l) = -\nu\\[4pt]
(f,\ l^3)^5 &= \frac{3}{3} (al)^2(al_1)^2(al_2)a_xl_x = \big((f,\ l^2)^4,\ l\big)\\
&= \Big(2n-\frac{2}{3}Bl+\frac{1}{3}Am,\ l\Big) = 2(n,\ l)+\frac{1}{3}A(m,\ l)\\
&= 2\mu-\frac{A}{3}\nu\,,
\end{aligned}
\]
und endlich

\[
\begin{aligned}
-(ak)(al)^2(al_1)^2(kl_2)^2a_xk_x &= \big((kl_2)^2k_x^2,\ (al)^2(al_1)^2a_x^2\big)\\
&= \big(m,\ (f,\ l^2)^4\big)\\
&= \Big(m,\ 2n-\frac{2B}{3}l+\frac{A}{3}m\Big) = 2(m,\ n)+\frac{2}{3}B(l,\ m)\\
&= 2\lambda+\frac{2}{3}B\nu\,.
\end{aligned}
\]

Das symbolische Product links ist aber ein Glied der Ueberschiebung
$\big((f,\ k),\ l^3\big)^6$, so dass damit der lineare Zusammenhang von $\lambda$
mit dieser Form des Systemes hergestellt ist.

Die sechs quadratischen Covarianten des Systemes werden also der Reihe nach
durch die Formen:

\[
\begin{aligned}
&l = (ak)^4a_x^2, \quad m = (kl)^2k_x^2, \quad n = (km)^2k_x^2 = (kl)^2(kk_1)^2k_{1x}^2\\
&\lambda = (m,\ n), \quad \mu = (n,\ l), \quad \nu = (l,\ m)
\end{aligned}
\]
ersetzt.


% ===== pdf idx 300, printed 286.0 =====
% !PRINTED 286
\weberpage{286}

264. \textit{Die fünf Invarianten.} Von den Formen

\[
(f,\ f)^6,\quad (k,\ k)^4,\quad \big((k,\ k)^2,\ k\big)^4,\quad (f,\ l^3)^6,\quad \big((f,\ k),\ l^4\big)^8
\]
werden, wie schon Eingangs erwähnt, die drei ersten beibehalten. Die vierte
ersetzen wir durch die Invariante $(l,\ n)^2 = A_{ln} = D$, die fünfte durch
die Combinante $-(lm)(mn)(nl) = R$. Was nämlich die Form $(f,\ l^3)^6$
betrifft, so hat man:

\[
\begin{aligned}
(f,\ l^3)^6 &= (al)^2(al_1)^2(al_2)^2 = \big((al)^2(al_1)^2a_x^2,\ l_{2x}^2\big)^2\\
&= \big((f,\ l^2)^4,\ l\big)^2 = \Big(2(\Delta,\ l)^2+\frac{A}{3}m,\ l\Big)^2\\
&= \Big(2n-\frac{2}{3}Bl+\frac{A}{3}m,\ l\Big)^2\\
&= 2(n,\ l)^2-\frac{2}{3}B(l,\ l)^2+\frac{1}{3}A(m,\ l)^2\,.
\end{aligned}
\]
Diese Gleichung lehrt, dass in der That an Stelle von $(f,\ l^3)^6$ die Form
$(n,\ l)^2 = A_{ln}$ treten darf.

In Bezug auf die Invariante $\big((f,\ k),\ l^4\big)^8$ ergiebt sich aber, dass
sie sich überhaupt nur um einen numerischen Factor von der Form
$R=-(lm)(mn)(nl)$ unterscheiden kann. Denn im Systeme von $f=a_x^6$ tritt nur
diese eine fundamentale schiefe Invariante auf, und sie ist vom Grade 15 in
den Coefficienten von $f$. Die Form $R$ ist aber gleichfalls schiefe Invariante
und vom Grade 15 in den Coefficienten; sie kann also nur eine lineare Function
der ersteren sein, d. h.

\[
R = \varrho\cdot\big((f,\ k),\ l^4\big)^8\,.
\]
Die fünf Invarianten des Systemes können demnach dargestellt werden durch die
symbolischen Producte

\[
\begin{aligned}
&A=(ab)^6,\quad B=(kk_1)^4\\
&C=(kk_1)^2(k_1k_2)^2(k_2k)^2,\quad D=A_{ln}\\
&R=-(lm)(mn)(nl)\,.
\end{aligned}
\]

265. \textit{Die zweiten Ueberschiebungen von $\Delta$ und $k$ über die
quadratischen Formen $l$, $m$, $n$.} Neben den fundamentalen quadratischen
Formen interessiren uns noch eine Reihe anderer Covarianten zweiten Grades,
die wir im Folgenden berechnen wollen. Es sind das in erster Linie die
zweiten Ueberschiebungen von $\Delta$ und $k$ über $l$, $m$, $n$, so weit sie
nicht bisher schon in Betracht gezogen wurden.

Was die Ueberschiebungen von $k$ betrifft, so war

\weq{1}{\begin{aligned}
&(k,\ l)^2 = m\,, \qquad (k,\ m)^2 = n\\[2pt]
&\text{Hiezu ergiebt sich}\qquad (k,\ n)^2 = \frac{Bm}{2}+\frac{Cl}{3}\,.
\end{aligned}}


% ===== pdf idx 301, printed 287.0 =====
% !PRINTED 287
\weberpage{287}

Denn es ist

\[
(k,\ n)^2 = (kn)^2k_x^2 = (kk_1)^2(k_1m)^2k_x^2 = (kk_1)^2(k_1k_2)^2(k_2l)^2k_x^2\,.
\]
Nach den Gesetzen der Reihenentwicklung hat man aber

\[
(kk_1)^2(k_1k_2)^2k_{2y}^2k_x^2 = \frac{1}{2}Bk_{y^2}+\frac{1}{3}C(yx)^2
\]
und daraus folgt unmittelbar für $y=l$ die obige Relation. Von den zweiten
Ueberschiebungen der Form $\Delta$ über $l$, $m$, $n$ ist $(\Delta,\ l)^2$
bereits berechnet [Nr. 263 (3)]. Es war

\[
(\Delta,\ l)^2 = n-\frac{B}{3}l\,.
\]
Hiezu findet sich:

\weq{2}{\begin{aligned}
(\Delta,\ m)^2 &= \frac{B}{6}m+\frac{C}{3}l\\
(\Delta,\ n)^2 &= \frac{B}{6}n+\frac{C}{3}m\,.
\end{aligned}}
Die zweite dieser Relationen erhalten wir, indem wir in der
Reihenentwicklung

\weq{$\alpha$}{(k\Delta)^2k_y^2\Delta_x^2 = \frac{B}{6}k_{y^2}+\frac{C}{3}(yx)^2}
$y$ durch $l$ und alsdann $(k\Delta)^2(kl)^2\Delta_x^2 = \big(\Delta,\
(kl)^2k_x^2\big)^2$ durch $(\Delta,\ m)^2$ ersetzen, während die dritte aus
derselben Entwicklung $(\alpha)$ für $y=m$ hervorgeht.

266. \textit{Die Ueberschiebungen von $f$ über Producte von $l$, $m$, $n$.}
Eine weitere Reihe quadratischer Covarianten liefern die vierten
Ueberschiebungen von $f$ über $l^2$, $lm$, $ln$, $m^2$, $mn$, $n^2$. Die drei
ersten ergeben sich unmittelbar durch zweimalige Ueberschiebung der schon
früher gegebenen Relation:

\[
(f,\ l)^2 = 2\Delta + \frac{A}{3}k
\]
über $l$, $m$ und $n$ unter Benutzung der eben gefundenen Relationen (1) und
(2). Das Resultat ist:

\weq{3}{\left.\begin{aligned}
(f,\ l^2)^4 &= -\frac{2}{3}Bl+\frac{A}{3}m+2n\\
(f,\ lm)^4 &= +\frac{2}{3}Cl+\frac{B}{3}m+\frac{A}{3}n\\
(f,\ ln)^4 &= +\frac{1}{9}ACl+\Big(\frac{2}{3}C+\frac{1}{6}AB\Big)m+\frac{B}{3}n\,.
\end{aligned}\right\}}
Hiezu treten durch nachfolgende Untersuchungen die drei weiteren Relationen:
%CONT>>


% ===== pdf idx 302, printed 288.0 =====
% !PRINTED 288
\weberpage{288}
%<<CONT

\weq{4}{\left.\begin{aligned}
(f,\ m^2)^4 &= \Big\{\frac{1}{2}A_{lm}+\frac{AC}{9}\Big\}l+\frac{1}{6}\Big\{\frac{AB}{3}-A_{ll}\Big\}m+\frac{B}{3}n\\
(f,\ mn)^4 &= \Big\{\frac{1}{2}A_{ln}+\frac{BC}{9}\Big\}l+\frac{1}{2}\Big\{\frac{B^2}{2}+\frac{AC}{3}\Big\}m+\frac{1}{3}\Big\{\frac{AB}{6}-A_{ll}\Big\}n\\
(f,\ n^2)^4 &= \Big\{\frac{1}{2}A_{mn}-\frac{C^2}{9}\Big\}l+\Big\{A_{ln}-\frac{CB}{18}\Big\}m+\frac{1}{3}\Big\{\frac{B^2}{6}-2A_{lm}\Big\}n
\end{aligned}\right\}}
wenn wir die Invarianten $(l,\ l)^2$, $(l,\ m)^2$, $(l,\ n)^2$ etc. resp. mit
$A_{ll}$, $A_{lm}$, $A_{ln}$ etc. bezeichnen.

Zur Ermittelung dieser drei letzten Ueberschiebungen ist es nöthig, zuvor
$(f,\ m)^2$ und $(f,\ n)^2$ zu berechnen. Wir gehen aus von der schon in
Nr. 255 benutzten Reihenentwicklung

\weq{$\beta$}{(ak)^3a_x^2k_xa_y = \frac{l}{4}(yx)\,.}
Ersetzen wir in ihr $y$ durch $l$, so erhält man:

\[
2(ak)^3a_x^2k_x(al)l_x = \frac{l^2}{2}\,.
\]
Entwickelt man die linke Seite nach dem Productsatz, so kommt:

\[
\begin{aligned}
\frac{l^2}{2} &= a_x^2(ak)^2\big\{(ak)^2l+(al)^2k_x^2-(kl)^2a_x^2\big\}\\
&= l^2+(ak)^2(al)^2a_x^2k_x^2-(ak)^2(kl)^2a_x^4\\
&= l^2+\big((f,\ l)^2,\ k\big)^2-\big(f,\ (kl)^2k_x^2\big)^2\\
&= l^2+\Big(2\Delta+\frac{A}{3}k,\ k\Big)^2-(f,\ m)^2\\
&= l^2+\frac{B}{3}\cdot k+\frac{A}{3}\Delta-(f,\ m)^2\,,
\end{aligned}
\]
daher

\weq{5}{(f,\ m)^2 = \frac{l^2}{2}+\frac{B}{3}\cdot k+\frac{A}{3}\Delta\,.}
Ersetzt man dagegen in $(\beta)$ $y$ durch $m$, so erhalten wir

\[
2(ak)^3(am)a_x^2k_xm_x = \frac{l\cdot m}{2}\,.
\]
Wendet man wiederum auf die linke Seite den Productsatz an, dann erhält man:

\[
\begin{aligned}
\frac{lm}{2} &= (ak)^2a_x^2\big\{(ak)^2m_x^2+(am)^2k_x^2-(km)^2a_x^2\big\}\\
&= (f,\ k)^4\cdot m + \big((f,\ m)^2,\ k\big)^2-\big(f,\ (km)^2k_x^2\big)^2\\
&= lm+\Big(\frac{l^2}{2}+\frac{B}{3}k+\frac{A}{3}\Delta,\ k\Big)^2-(f,\ n)^2\,.
\end{aligned}
\]
Im mittleren Gliede rechts ist die Ueberschiebung $(l^2,\ k)^2$ zu berechnen.


% ===== pdf idx 303, printed 289.0 =====
% !PRINTED 289
\weberpage{289}

Wir benutzen hiezu die Reihe:
\weq{II}{l_x l_y \cdot l_x l_y = (l^2)_{y^2} + \frac{1}{3} A_{ll} (yx)^2.}

Ersetzt man hier $y$ durch $k$, so kommt

\[
(lk)^2 k_x^2 \cdot l_x^2 = (l^2,\, k)^2 + \frac{1}{3}\left(2C+\frac{AB}{3}\right)k_x^4,
\]

oder

\[
(l^2,\, k)^2 = lm - \frac{1}{3}\left(2C+\frac{AB}{3}\right)k.
\]

Trägt man diesen Werth in (6) ein, so erhält man:
\[
\frac{lm}{2} = lm + \left\{\frac{1}{2}\left(lm-\frac{2}{3}Ck-\frac{1}{9}ABk\right)+\frac{B}{3}\Delta+\frac{AB}{18}k\right\}-(f,\,n)^2.
\]
Daraus folgt endlich:
\weq{7}{(f,\,n)^2 = lm - \frac{1}{3}Ck + \frac{B}{3}\Delta.}

Mit Hilfe der Relationen (5) und (7) lassen sich nun in einfacher Weise die Ueberschiebungen $(f,\,m^2)^4$, $(f,\,mn)^4$, $(f,\,n^2)^4$ berechnen.

Es ist:
\[
(f,\,m^2)^4 = \frac{1}{2}(l^2,\,m)^2+\frac{B}{3}(k,\,m)^2+\frac{A}{3}(\Delta,\,m)^2.
\]

Hierin ist, wenn wir in der Reihenentwicklung (II) $y$ durch $m$ ersetzen:
\[
(l^2,\,m)^2=(l,\,m)^2\cdot l - \frac{1}{3}A_{ll}m.
\]

Ferner hatten wir in (2) Nr. 265
\[
(\Delta,\,m)^2=\frac{B}{6}m+\frac{C}{3}l,
\]
und da ausserdem $(k,\,m)^2=n$, $(l,\,m)^2=A_{lm}$, so ist endlich:
\[
(f,\,m^2)^4=\frac{1}{2}\left\{A_{lm}\cdot l - \frac{1}{3}A_{ll}\cdot m\right\}+\frac{B}{3}n+\frac{A}{3}\left\{\frac{B}{6}m+\frac{C}{3}l\right\}.
\]

Schieben wir aber die Relation (5) zweimal über $n$, so kommt:
\[
(f,\,m\cdot n)^4=\frac{1}{2}(l^2,\,n)+\frac{B}{3}(k,\,n)^2+\frac{A}{3}(\Delta,\,n)^2.
\]

Aus der Reihe (II) folgt wiederum für $y=n$
\[
(l^2,\,n)=(l,\,n)^2\cdot l - \frac{1}{3}A_{ll}n.
\]

Die Ueberschiebung wird also unter Berücksichtigung der Relationen (1) und (2) Nr. 265
\[
(f,\,mn)^4=\frac{1}{2}\left\{A_{ln}l-\frac{1}{3}A_{ll}n\right\}+\frac{B}{3}\left(\frac{B}{2}m+\frac{C}{3}l\right)+\frac{A}{3}\left(\frac{B}{6}n+\frac{C}{3}m\right)\cdot
\]

Endlich ergiebt sich aus der Gleichung (7):
\[
(f,\,n^2)^2=(lm,\,n)^2-\frac{1}{3}C(k,\,n)^2+\frac{B}{3}(\Delta,\,n)^2.
\]


% ===== pdf idx 304, printed 290.0 =====
% !PRINTED 290
\weberpage{290}

Zur Berechnung des ersten Gliedes der rechten Seite benutzen wir die Reihe
\[
lm_y^2+ml_y^2=2(lm)_{y^2}+\frac{2}{3}A_{lm}(yx)^2.
\]

Ersetzt man hier $y$ durch $n$, so kommt:
\[
2(lm,\,n)^2=lA_{mn}+mA_{ln}-\frac{2}{3}A_{lm}n.
\]

Demnach wird
\[
(f,\,n^2)^2=\frac{1}{2}\left\{A_{mn}l+A_{ln}m-\frac{2}{3}A_{lm}n\right\}-\frac{C}{3}\left\{\frac{C}{3}l+\frac{B}{2}m\right\}+\frac{B}{3}\left\{\frac{B}{6}n+\frac{C}{3}m\right\}.
\]

Damit wollen wir die Berechnung von quadratischen Covarianten der Form sechster Ordnung abschliessen, und zur Berechnung der Invarianten der quadratischen Fundamentalcovarianten übergehen.

267. \textit{Berechnung der Invarianten des simultanen Systemes $l$, $m$, $n$.} Die drei quadratischen Covarianten besitzen ausser der schiefen Invariante $R=(lm)(mn)(ln)$ noch sechs simultane Invarianten (vgl. Nr. 131)
\[
A_{ll},\quad A_{lm},\quad A_{ln},\quad A_{mm},\quad A_{mn},\quad A_{nn},
\]
die wir nun im Folgenden durch die Fundamentalinvarianten des Systemes von $f=a_x^6$ darstellen wollen.

Wir hatten bereits früher erhalten (vgl. Nr. 258 und Nr. 264)
\weq{1}{\left.\begin{aligned} A_{ll}&=2C+\frac{1}{3}AB\\ A_{lm}&=(k,\,l^2)^4=\frac{2}{3}(B^2+AC)\\ A_{ln}&=D. \end{aligned}\right\}}

Hiezu ergiebt sich:
\weq{2}{\left.\begin{aligned} A_{mm}&=D\\ A_{mn}&=\frac{1}{3}\left\{B^3+2C^2+\frac{4}{3}ABC\right\}\\ A_{nn}&=\frac{1}{2}BD+\frac{2}{9}C(B^2+AC) \end{aligned}\right\}}

Denn man hat:
\[
A_{mm}=((k,\,l)^2,\,m)^2=(k,\,ml)^4=((k,\,m)^2,\,l)^2=(n,\,l)^2
\]
\[
A_{mn}=((k,\,l)^2,\,n)^2=(k,\,nl)^4=((k,\,n)^2,\,l)^2=\left(\frac{B}{2}m+\frac{C}{3}l,\ l\right)^2
\]
\[
A_{nn}=((k,\,m)^2,\,n)^2=(k,\,mn)^4=((k,\,n)^2,\,m)^2=\left(\frac{B}{2}m+\frac{C}{3}l,\ l\right)^2.
\]


% ===== pdf idx 305, printed 291.0 =====
% !PRINTED 291
\weberpage{291}

268. \textit{Die Relationen für die schiefe Invariante.} Aus der Theorie der simultanen quadratischen Formen ergiebt sich unmittelbar, dass (vgl. Nr. 133) $R$ auch dargestellt werden kann durch
\[
R=A_{l\lambda}=A_{m\mu}=A_{n\nu},
\]
und dass ferner zwischen $R^2$ und den übrigen fundamentalen Covarianten eine Relation bestehen muss, welche aus
\[
2R^2=\begin{vmatrix} A_{ll}, & A_{lm}, & A_{ln} \\ A_{ml}, & A_{mm}, & A_{mn} \\ A_{nl}, & A_{nm}, & A_{nn} \end{vmatrix}
\]
durch Substitution der in Nr. 267 berechneten Werthe der Determinantenelemente hervorgeht.

Weitere Relationen will ich nicht mehr anführen. Die schon früher in der Theorie der quadratischen Formen gewonnenen zahlreichen Beziehungen lassen sich hier direct übertragen. In ähnlicher Weise würde man auch die Theorie der biquadratischen Formen benutzen können, um Relationen zwischen den Covarianten vierten Grades des Systemes aufzustellen.

\S\,28. \textbf{Die Discriminante der Form sechsten Grades und irrationale Relationen.}

269. \textit{Vorbereitende Untersuchungen.} Die erste Berechnung der Discriminante einer Form sechsten Grades haben wir von Brioschi. Er benutzte hiebei die allgemeinen Differentialgleichungen, denen Invarianten überhaupt genügen, und charakterisirte dann die Discriminante, indem er ihnen specielle hinzufügte, welche nur durch Discriminanten befriedigt werden (Crelle's Journal Bd. 53. Vgl. auch Faà di Bruno, \glqq Theorie der binären Formen\grqq{}, deutsche Uebersetzung von Walter, S. 275).

Die genaue Kenntniss des vollen Systemes von $f$ gestattet uns indess hier wiederum den nun schon mehrmals eingeschlagenen directen Weg zu verfolgen (vgl. Discrim. der Form vierten und fünften Grades), indem wir uns fragen:

\begin{quote}
\glqq Welche Relation besteht zwischen den Invarianten $A$, $B$, $C$, $D$, wenn $f=\alpha_x^2\cdot\varphi$, wo $\alpha_x$ eine lineare und $\varphi$ eine biquadratische Form ist?\grqq{}
\end{quote}

Wie bei der Form fünften Grades, und wie überhaupt allgemein (vgl. Nr. 99), gehen aus der Voraussetzung $f=\alpha_x^2\cdot\varphi$ folgende vier wichtige Beziehungen hervor:
%CONT>>


% ===== pdf idx 306, printed 292.0 =====
% !PRINTED 292
\weberpage{292}
%<<CONT

\weq{1}{(a\alpha)^6=0,\quad (a\alpha)^5a_x=0,}
\weq{2}{(a\alpha)^4a_x^2=\lambda\cdot\alpha_x^2,}
\weq{3}{\frac{5}{2}(ab)^2(a\alpha)^4b_x^4 - \frac{5}{2}(k\alpha)^2k_x^2\cdot\alpha_x^2=\frac{1}{3}A\cdot\alpha_x^4,}

wobei Gleichung (3) aus $(a\alpha)^4a_x$ durch die in Nr. 99 angegebene Reihenentwicklung gewonnen wird. Der Werth des Parameters $\lambda$ berechnet sich gleichfalls wie früher. Wir ersetzen in (2) $x$ durch $b$ und multipliciren mit $b_x^4$, dann erhalten wir:
\[
(a\alpha)^4(ab)^2b_x^4=\lambda\cdot(f,\,\alpha^2)^2
\]
und demnach wegen (3)
\weq{4}{\lambda\cdot(f,\,\alpha^2)^2=(k\alpha)^2k_x^2\cdot\alpha_x^2+\frac{2}{15}A\cdot\alpha_x^4.}

Ueberschiebt man aber diese Relation zweimal über $\alpha_x^2$, so kommt in Verbindung mit Gleichung (2) direct:
\weq{5}{\lambda^2=\frac{1}{6}(k\alpha)^4.}

Die nächste Aufgabe wäre nun, wiederum auch eine Relation für die erste Ueberschiebung $(f,\,\alpha)$ zu ermitteln und auf Grund der so berechneten Ueberschiebungen von $f$ und $\alpha$ das Formensystem der simultanen Invarianten von $f$ und $\alpha$ mit den Invarianten $A$, $B$, $C$, $D$ in Beziehung zu setzen. Hiebei zeigt sich aber, dass insbesondere die höheren Ueberschiebungen von $k$, $\Delta$ und $l$ über $\alpha$ die hervorragendste Rolle spielen, und wir suchen demnach hier von vornherein sie mit dem allgemeinen Formensystem in Beziehung zu setzen.

270. \textit{Berechnung von $(l\alpha)^2$ und $(\Delta\alpha)^2$.} Zur Berechnung der Ueberschiebungen von $l$ über $\alpha$ benutzen wir am geschicktesten die Relation (4), indem wir sie dreimal über $k$ schieben. Dann kommt:
\weq{1}{\begin{aligned} \lambda(a\alpha)^2(ak)^3a_xk_x&=\frac{1}{2}\left\{(kk_1)^2(\alpha k_1)\alpha_x+(kk_1)(\alpha k_1)^2k_x\right\}(k\alpha)^2k_{1x}-\frac{2}{15}A(k,\,\alpha^3)^3\alpha_x\\ &=-\frac{1}{2}(\Delta,\,\alpha^3)^3\alpha_x-\frac{2}{15}A(k,\,\alpha^3)^3\alpha_x,\end{aligned}}

weil das zweite Glied der Klammer rechts identisch verschwindet und $\Delta=(kk_1)^2k_x^2k_{1x}^2$ ist. Anderntheils ist aber, da nach Nr. 253, $(f,\,k)^3=0$,
\[
(ak)^3a_y^2k_xa_x=-\frac{1}{2}l_xl_y\cdot(xy),
\]
und demnach
\weq{2}{(a\alpha)^2(ak)^3a_xk_x=\frac{1}{2}(l,\,\alpha)\cdot\alpha_x.}

Führt man diesen Werth in (1) ein, so kommt nach Division mit $\frac{1}{2}\alpha_x$:
\weq{3}{\lambda\cdot(l,\,\alpha)=-\left\{(\Delta,\,\alpha^3)^3+\frac{4}{15}A\cdot(k,\,\alpha^3)^3\right\}}


% ===== pdf idx 307, printed 293.0 =====
% !PRINTED 293
\weberpage{293}

und hieraus durch abermalige Ueberschiebung mit $\alpha$:
\weq{4}{\lambda\cdot(l,\ \alpha^2)^2=-\left\{(\varDelta,\ \alpha^4)^4+\frac{4}{15}A\cdot(k,\ \alpha^4)^4\right\}\cdot}
Eine zweite Gleichung zwischen $(l\alpha)^2$ und $(\varDelta\alpha)^4$ erhalten wir durch
viermalige Ueberschiebung von Relation (4) mit der Form $k$, nämlich
\weq{5}{\lambda\cdot(l\alpha)^2=(\varDelta,\ \alpha^4)^4+\frac{2}{15}A\cdot(k,\ \alpha^4)^4.}
Indem wir beide Gleichungen addiren, erhalten wir wegen $(k\alpha)^4=6\lambda^2$
\weq{6}{(l\alpha)^2=-\frac{2}{5}A\cdot\lambda,}
und indem wir sie subtrahiren
\[
(\varDelta\alpha)^4=-\frac{6}{5}A\cdot\lambda^2.
\]
Nun lässt sich aber Gleichung (3) auch in der Form darstellen
\[
\left(\lambda\cdot l+(\varDelta,\ \alpha^2)^2+\frac{4}{15}A\cdot(k,\ \alpha^2)^2,\ \alpha\right)=0,
\]
d.\,h. (vgl. Nr. 99) die quadratische Form
\[
\lambda\cdot l_x^2+(\varDelta\alpha)^2\varDelta_x^2+\frac{4A}{15}(k\alpha)^2k_x^2
\]
ist eine Potenz von $\alpha$, so dass also, wenn $\varrho$ eine Proportionalitätsconstante
, die Relation besteht
\weq{7}{\varrho\cdot\alpha_x^2-\lambda\cdot l_x^2=(\varDelta\alpha)^2\varDelta_x^2+\frac{4}{15}A\cdot(k\alpha)^2k_x^2.}
Die Constante $\varrho$ finden wir sofort, wenn wir diese Relation (7) zweimal
über sich selbst schieben. Hiebei ergiebt sich links, da
\[
(l,\ l)^2=2C+\frac{1}{3}AB,
\]
der Werth:
\weq{8}{-2\lambda\cdot\varrho(l\alpha)^2+\lambda^2\left(2C+\frac{1}{3}AB\right);}
rechts dagegen erhalten wir:
\[
(\varDelta\alpha)^2(\varDelta_1\alpha)^2(\varDelta\varDelta_1)^2+\frac{8}{15}A\cdot(\varDelta\alpha)^2(k\alpha)^2(\varDelta k)^2+\frac{16}{225}A^2\cdot(k\alpha)^2(k_1\alpha)^2(kk_1)^2,
\]
oder, weil nach der Theorie der biquadratischen Formen $(\varDelta,\ k)^2=\frac{B}{6}k$
und $(\varDelta,\ \varDelta)^2=\frac{C}{3}k-\frac{B}{6}\varDelta$:
\weq{9}{\frac{C}{3}(k\alpha)^4-\frac{B}{6}(\varDelta\alpha)^4+\frac{8}{90}AB\cdot(k\alpha)^4+\frac{16}{225}A^2\cdot(\varDelta\alpha)^4.}
Tragen wir nun in (8) und (9) die Werthe von $(k\alpha)^4$, $(\varDelta\alpha)^4$, $(l\alpha)^2$
ein und setzen beide Ausdrücke wieder einander gleich, so kommt
nach einfacher Reduction:
\weq{10}{\varrho=\frac{1}{2}B-\frac{8}{75}A^2.}


% ===== pdf idx 308, printed 294.0 =====
% !PRINTED 294
\weberpage{294}

271. \textit{Die Discriminante.} Die Relation (7) gestattet nun aber
auch die Invariante $D=(mm_1)^2$ mit den simultanen Invarianten
$(\varDelta\alpha)^4$ und $(k\alpha)^4$ in Beziehung zu setzen und damit erreichen wir zugleich
 das Ziel unserer Aufgabe. Ueberschiebt man nämlich die erwähnte
 Relation (7) zweimal über $k$, so erhält man wegen $(k,\ l)^2=m$
zunächst:
\weq{1}{\varrho\cdot(k,\ \alpha^2)^2-\lambda\cdot m=(\varDelta\alpha)^2(\varDelta k)^2k_x^2+\frac{4}{15}A\cdot(k\alpha)^2(k_1k)^2k_x^2.}
Die Werthe der beiden Terme rechts ergeben sich aus den entsprechenden
 Reihenentwicklungen:
\[
\begin{aligned}
(\varDelta k)^2\varDelta_y^2k_x^2 &= \{(\varDelta,\ k)^2\}_{y^2}+\frac{1}{3}(\varDelta,\ k)^4(yx)^2=\frac{B}{6}k_x^2k_y^2+\frac{C}{3}(yx)^2\\
(kk_1)^2k_x^2k_{1y}^2 &= \{(k,\ k)^2\}_{y^2}+\frac{1}{3}(k,\ k)^4(xy)^2=\varDelta_x^2\varDelta_y^2+\frac{B}{3}(yx)^2,
\end{aligned}
\]
indem man darin $y$ durch $\alpha$ ersetzt. Führt man alsdann diese
Werthe in (1) ein und ersetzt $\varrho$ durch $\frac{1}{2}B-\frac{8}{75}A^2$, so kommt
nach Multiplication mit $-3$:
\[
3\lambda m=\left(B-\frac{8}{25}A^2\right)(k\alpha)^2k_x^2-\frac{4}{5}A\cdot(\varDelta\alpha)^2\varDelta_x^2-\left(C+\frac{4}{15}AB\right)\alpha_x^2.
\]
Bilden wir nun auf beiden Seiten dieser Gleichung
\[
3\lambda m=P\cdot(k\alpha)^2k_x^2+Q\cdot(\varDelta\alpha)^2\varDelta_x^2+R\cdot\alpha_x^2
\]
die Hesse'sche Form, so erhalten wir
\[
\begin{aligned}
9\lambda^2D &= P^2(\varDelta\alpha)^4+Q^2((\varDelta,\ \varDelta)^2,\ \alpha^4)^4\\
&\quad+2P\cdot Q\cdot((\varDelta,\ k)^2,\ \alpha^4)^4+2P\cdot R\cdot(k\alpha)^4+2QR(\varDelta\alpha)^4.
\end{aligned}
\]
Ersetzen wir hierin $(\varDelta,\ \varDelta)^2$ und $(\varDelta,\ k)^2$ durch ihre Werthe
$\frac{C}{3}k-\frac{B}{6}\varDelta$, resp. $\frac{B}{6}k$ und ordnen alsdann nach den Ueberschiebungen
$(\varDelta\alpha)^4$ und $(k\alpha)^4$, so geht diese Gleichung über in:
\[
9\lambda^2D=(\varDelta\alpha)^4\left\{P^2-\frac{B}{6}Q^2+2QR\right\}+(k\alpha)^4\left\{\frac{C}{3}\cdot Q^2+\frac{B}{3}\cdot PQ+2PR\right\}
\]
oder wenn wir nun $(\varDelta\alpha)^2$, $(k\alpha)^2$, $P$, $Q$, $R$ durch ihre Werthe ersetzen
und mit $\lambda^2$ dividiren: % checked against the GDZ scan (27 Oct 2026): the only mark over an equals sign on this page, a blob above the second "=" of the (\Delta k)^2 line, is an ink speck (that equality is exact, and the printed \doteq of this book has a short slanted stroke at the left); plain "="
\[
\begin{aligned}
9D &= -\frac{6}{5}A\left\{B^2-\frac{8}{25}\cdot A^2B+\frac{64}{625}A^4+\frac{8}{5}AC\right\}\\
&\quad+6\left\{\frac{4}{3}\cdot\frac{16}{25}A^2C-\frac{4}{5}AB^2+\frac{32}{125}A^3B-2BC\right\}.
\end{aligned}
\]
Das ist die Invariantenrelation, welche die Voraussetzung $f=\alpha_x^2\cdot\varphi$
zur Folge hat. Wir können sie auf die einfachere Form bringen,
indem wir $\frac{2}{5}A=A_1$ setzen:
%CONT>>


% ===== pdf idx 309, printed 295.0 =====
% !PRINTED 295
\weberpage{295}
%<<CONT

\[
12A_1^5-30A_1^3B-20A_1^2C+15A_1B^2+12BC+9D=0,
\]
und die linke Seite dieser Gleichung müssen wir als Discriminante
von $f$ bezeichnen, mit der sie auch in Bezug auf den Grad in den
Coefficienten übereinstimmt.

272. \textit{Irrationale Relationen bei Combinanten.} Ehe wir die Untersuchungen
 über Co- und Invariantenrelationen verlassen, will ich
noch auf eine Gattung von irrationalen Beziehungen zwischen Combinanten
 hinweisen, auf welche bereits \gesperrt{Brill}, \gesperrt{Stephanos}, \gesperrt{Meyer}
und \gesperrt{Stroh} aufmerksam gemacht haben.

Sind nämlich $f=r_x^n$ und $\varphi=s_x^n$ zwei binäre Formen $n^{\mathrm{ten}}$ Grades,
so ist nach einem Satze von \gesperrt{Gordan} (vgl. Nr. 68 und Math. Ann.
Bd. V) jede Combinante der beiden Formen Covariante des Cayley'schen
Ausdruckes
\[
\varPsi=f(x)\varphi(y)-\varphi(x)f(y),
\]
also seiner Elementarcovarianten. Diese Elementarcovarianten sind
in unserem Falle die ungeraden Ueberschiebungen von $f$ über $\varphi$, da
die geraden sich gegenseitig aufheben, also die Formen:
\[
(f,\ \varphi),\quad (f,\ \varphi)^3,\quad (f,\ \varphi)^5\ \ldots\ \text{etc.}
\]
Aber diese Covarianten können nicht von einander unabhängig sein.
Denn wenn $\varphi$ proportional $f$ ist, so verschwindet $\varPsi$ und demnach alle
Elementarcovarianten, und umgekehrt: Ist $\varPsi$ identisch null, so ist
$f=\varrho\cdot\varphi$, wo $\varrho$ eine Constante. Damit aber $f$ und $\varphi$ einander proportional
 sind, ist es bekanntlich schon hinreichend, dass ihre Functionaldeterminante
 $(f,\ \varphi)$ identisch verschwindet. Wenn also schon das
Verschwinden der ersten Elementarcovariante auch das Verschwinden
von $\varPsi$ und damit aller übrigen Elementarcovarianten nach sich zieht,
so können diese letzteren nicht mehr unabhängig sein von $(f,\ \varphi)$.
Es müssen vielmehr die Coefficienten der späteren Ueberschiebungen
homogene und irrationale Functionen der Coefficienten von $(f,\ \varphi)$ sein.

Man kann daher sich die Fragen stellen: Welches sind diese
Coefficientenrelationen und wie gelangt man zu ihnen? Was die zweite
Frage betrifft, so kann man zu ihrer Beantwortung einen schon mehrfach
 benutzten Weg einschlagen. Wir bilden eine Covariante von
$(f,\ \varphi)$ durch Ueberschiebung dieser Form über sich selbst, und sehen
nach, ob sich diese Form nicht auch simultan durch eine Ueberschiebung
 von $F=(f,\ \varphi)$ über eine der übrigen Elementarcovarianten
darstellen lässt. Dies ist sicher der Fall, wenn z.\,B. ein Glied der
Ueberschiebung $(F,\ F)^\nu$ verschwindet, da man weiss, dass jedes Glied
einer Ueberschiebung nach dieser und niedrigeren Ueberschiebungen
von $F$ über andere Formen entwickelt werden kann.


% ===== pdf idx 310, printed 296.0 =====
% !PRINTED 296
\weberpage{296}

Nun tritt bereits für $\nu=4$, also in $(F,\ F)^4$ ein solches Glied
auf. Es ist, wenn $F=(f,\ \varphi)=(rs)r_x^{n-1}s_x^{n-1}$, das Glied
\[
(rs)(r_1s_1)(rr_1)(sr_1)(rs_1)(ss_1)r_x^{n-3}s_x^{n-3}r_{1x}^{n-3}s_{1x}^{n-3},
\]
welches verschwindet, weil es durch Vertauschung von $r$ mit $r_1$ nur
sein Zeichen ändert. Es muss also eine Relation bestehen von der
Form
\weq{I}{0=(F,\ F)^4+\sum_{\nu=0}^{\nu=3}\varrho_\nu(F,\ \varPhi_\mu)^\nu,}
wobei $\varPhi_\mu$ irgend welche spätere Elementarcovarianten niedrigerern % [sic]
Grades als $F$ sind. Durch Coefficientenvergleichung findet man alsdann
 jene Relationen, welche die Coefficienten der Formen $\varPhi_\mu$ mit
denen von $F$ in Beziehung setzen.

273. \textit{Beispiel.} Wir wollen hier an einem speciellen Falle, in
welchem die Functionaldeterminante $F=(f,\ \varphi)$ eine Form sechsten
Grades ist, diese allgemeinen Betrachtungen erläutern. Die Formen
$f$ und $\varphi$ sind alsdann zwei biquadratische Formen
\[
f=r_x^4,\quad \varphi=s_x^4,
\]
und ihre Combinanten entstehen durch Ueberschiebung der beiden
einzig vorhandenen Elementarcombinanten
\[
\begin{aligned}
F &= (rs)r_x^3s_x^3=a_x^6\\
\varPhi &= (rs)^3r_xs_x.
\end{aligned}
\]
Nach dem Vorausgegangenen sind die Coefficienten der Form $\varPhi$
irrationale Functionen der Coefficienten von $F$, und unsere Aufgabe
ist, die Beziehungen zwischen diesen Coefficienten herzustellen. Da
wir es hier ausser der Combinante $F$ nur noch mit einer Combinante
$\varPhi$ zu thun haben, so wird die oben aufgestellte allgemeine Relation
(I) in einfacher Weise zum Ziele führen. Wir bilden demnach $(F,\ F)^4$,
also die Covariante $k$ von $F=a_x^6$. Da diese Form $k=(ab)^4a_x^2b_x^2$
aus der gemischten Polare $a_x^2a_y^2a_z^2$ für $y=z=b$ hervorgeht, so wollen
wir zunächst diese gemischte Polare in Symbolen $r$, $s$ darstellen.

Wir substituiren zu dem Zwecke in der Identität:
\[
a_x^6=(rs)r_x^3s_x^3
\]
für $x_i$ den Werth $\lambda_1x_i+\lambda_2y_i+\lambda_3z_i$, dann kommt:
\[
(\lambda_1a_x+\lambda_2a_y+\lambda_3a_z)^6=(rs)(\lambda_1r_x+\lambda_2r_y+\lambda_3r_z)^3(\lambda_1s_x+\lambda_2s_y+\lambda_3s_z)^3.
\]
Durch Comparation der Glieder mit dem Factor $\lambda_1^2\lambda_2^2\lambda_3^2$ erhält man:
\[
90a_x^2a_y^2a_z^2=36(rs)r_xr_yr_zs_xs_ys_z+9(rs)\sum r_x^2r_ys_ys_z^2,
\]
%CONT>>


% ===== pdf idx 311, printed 297.0 =====
% !PRINTED 297
\weberpage{297}
%<<CONT
wobei sich die Summe über die sechs möglichen Producte von der Gestalt $r_x^2 r_y s_y s_z^2$ erstreckt. Addiren wir rechts sechsmal das Product $9\,(rs)\,r_x r_y r_z s_x s_y s_z$, indem wir gleichzeitig jeden solchen Summanden von je einem Gliede unter dem Summenzeichen subtrahiren, so findet man nach zweimaliger Anwendung der Identitäten

\[
r_x s_y - s_x r_y = (rs)(xy), \quad r_x s_z - s_x r_z = (rs)(xz), \quad r_y s_z - r_z s_y = (rs)(yz)
\]
die Relation:
\[
90\,a_x^6\,a_y^2\,a_z^2 = 90\,(rs)\,r_x r_y r_z s_x s_y s_z + 9\,(rs)^3 \sum (xy)^2 r_z s_z.
\]
Dividiren wir dieselbe durch 45 und setzen nun
\[
\frac{2}{5}\,\Phi = \frac{2}{5}\,(rs)^3\,r_x s_x = p_x^2\text{, etc.},
\]
so kommt:
\[
2\,a_x^2 a_y^2 a_z^2 = 2\,(rs)\,r_x r_y r_z s_x s_y s_z + \frac{1}{2}\left\{ p_x^2 (yz)^2 + p_y^2 (xz)^2 + p_z^2 (xy)^2 \right\},
\]
oder, indem wir den Productsatz
\[
p_y p_z (xy)(xz) + \frac{1}{2} p_x^2 (yz)^2 = \frac{1}{2}\left\{ p_y^2(xz)^2 + p_z^2(xy)^2 \right\}
\]
auf den zweiten Term rechts anwenden:
\weq{II}{
2\,a_x^2 a_y^2 a_z^2 = 2\,(rs)\,r_x r_y r_z s_x s_y s_z + p_x^2 (yz)^2 + p_y p_z (xy)(xz).
}
Ersetzen wir hierin $y$ und $z$ durch $b$ und multipliciren mit $b_x^2$, so kommt:
\weq{III}{
2\,(ab)^4 a_x^2 b_x^2 = 2\,(rs)(rb)^2(sb)^2 r_x s_x b_x^2 + (bp)^2 b_x^4.
}
Der zweite Term rechts ist nichts anderes als $\frac{2}{5}(F,\Phi)^2$. Der Werth des ersten Termes ergiebt sich aus derselben Gleichung (II), wenn wir darin $y$ durch $r$, $z$ durch $s$ ersetzen und mit $(rs)\,r_x s_x$ multipliciren; man erhält alsdann:
\weq{1}{
2\,(rs)(ar)^2(as)^2 a_x^2 r_x s_x = p_x^2\cdot(rs)^3 r_x s_x + (pr)(ps)(rs)\,r_x^2 s_x^2.
}
Hierin ist der zweite Term rechts ein Glied der Ueberschiebung
\[
(f,p)^2 = (rs)(r_x^3 s_x^3, p_x^2)^3 = \frac{(rs)}{5}\left\{ (pr)^2 r_x s_x^3 + (ps)^2 s_x r_x^3 + 3(pr)(ps) r_x^2 s_x^2 \right\}.
\]
Ersetzt man in der Klammer rechts die Grösse $(pr)^2 s_x^2 + (ps)^2 r_x^2$ nach dem Productsatz durch $2(pr)(ps) r_x s_x + (rs)^2 p_x^2$, so kommt:
\weq{2}{
(f,p)^2 = (ps)(pr)(rs)\,r_x^2 s_x^2 + \frac{1}{5}(rs)^3\,r_x s_x\cdot p_x^2.
}
Führt man also in die Gleichung (1) den aus dieser Gleichung (2) sich ergebenden Werth von $(ps)(pr)(rs)\,s_x^2 r_x^2$ ein, so erhalten wir
\[
2\,(rs)(ar)^2(as)^2\,a_x^2 r_x s_x = \frac{4}{5}\,p_x^2\cdot(rs)^3\,r_x s_x + (f,p)^2,
\]
%CONT>>


% ===== pdf idx 312, printed 298.0 =====
% !PRINTED 298
\weberpage{298}
%<<CONT
oder, weil $p_x^2 = p = \frac{2}{5}(rs)^3 r_x s_x$,
\[
2\,(rs)(ar)^2(as)^2\,a_x^2 r_x s_x = 2p^2 + (f,p)^2 = \frac{8}{25}\,\Phi^2 + \frac{2}{5}(F,\Phi)^2.
\]
Durch Substitution dieses Werthes in (III) ergiebt sich endlich nach Division mit dem Factor 2:
\weq{IV}{
(F,F)^4 = \frac{4}{25}\,\Phi^2 + \frac{2}{5}(F,\Phi)^2,
}
und dies ist die gesuchte Fundamentalrelation zwischen den Coefficienten von $F$ und $\Phi$.

274. Führen wir in dieselbe die Symbole $a$, $k$ und $p$ von $F$ resp. $(F,F)^4$ und $\frac{2}{5}\Phi$ ein, so nimmt sie die Gestalt an
\weq{3}{
k = p^2 + (ap)^2\,a_x^4.
}
Durch Vergleichung der Coefficienten von $x$ auf beiden Seiten dieser Identität erhält man fünf Gleichungen zur Berechnung der drei Grössen $p_0\,p_1\,p_2$. Doch lässt sich dieses System in einfacher Weise auf eine einzige Gleichung für die Discriminante $P=(pp_1)^2$ der quadratischen Form $p$ reduciren. Ich überschiebe diese Identität zunächst viermal über $f$, und erhalte, da $(k,f)^4=(ak)^4 a_x^2=l$,
\[
l = (f,p^2)^4 + (f,(f,p)^2)^4 = (ap)^2(ap_1)^2 a_x^2 + (ab)^4(bp)^2 a_x^2.
\]
Der Werth des ersten Termes rechts folgt aus der Fundamentalrelation (3), indem man sie zweimal über $p$ schiebt:
\[
(f,p^2)^4 = (k,p)^2 - \frac{2}{3}\,P\cdot p.
\]
Um das zweite Glied zu berechnen, bilde ich die Ueberschiebung
\[
\begin{aligned}
(k,p)^2 &= (ab)^4(a_x^2 b_x^2,\, p_x^2) = \frac{(ab)^4}{3}\left\{ (bp)^2 a_x^2 + 2(ap)(bp)\,a_x b_x \right\} \\
&= \frac{(ab)^4}{3}\left\{ (bp)^2 a_x^2 + (ap)^2 b_x^2 + (bp)^2 a_x^2 - (ab)^2 p \right\}.
\end{aligned}
\]
Hieraus folgt direct:
\[
(ab)^4(bp)^2 a_x^2 = (f,(f,p)^2)^4 = (k,p)^2 + \frac{A}{3}\,p.
\]
Demnach wird
\weq{V}{
l = 2(k,p)^2 + \frac{1}{3}\,p\,(A - 2P).
}
Bildet man nun die zweite Ueberschiebung dieser Relation über $p$, so findet man:
\weq{4}{
(lp)^2 = 2(k,p^2)^4 + \frac{P}{3}(A-2P).
}
Die Fundamentalrelation (3) liefert aber für $(k,p^2)^4$ unmittelbar durch viermalige Ueberschiebung über $k$:
\[
(k,p^2)^4 = B - (lp)^2.
\]


% ===== pdf idx 313, printed 299.0 =====
% !PRINTED 299
\weberpage{299}
Trägt man diesen Werth in die letzte Gleichung (4) ein, so kommt:
\weq{VI}{
3(lp)^2 = 2B + \frac{P}{3}(A-2P).
}
Gelingt es uns nun, die Ueberschiebung $(lp)^2$ noch auf eine zweite Art darzustellen, so erhalten wir in Verbindung mit Gleichung (VI) eine Relation für $P$.

275. Nun erinnern wir uns aber daran, dass die zweimalige Ueberschiebung von $\Delta$ über $l$ dargestellt ist durch (vergl. Nr. 265, 2)
\[
(\Delta,l)^2 = n - \frac{B}{3}\,l,
\]
dass also rechts wiederum die Covariante $l$ auftritt.

Wenn wir daher die Relation (V) zweimal über $\Delta$ schieben, so erhalten wir eine neue Relation für $l$, nämlich:
\weq{5}{
n - \frac{1}{3}\,Bl = 2\big((k,p)^2,\Delta\big)^2 + \frac{1}{3}(\Delta,p)^2(A-2P).
}
Es ist aber:
\[
\big((k,p)^2,\Delta_x^4\big)^2 = (kp)^2(k\Delta)^2\Delta_x^2.
\]
Der Term rechts ist ein Glied der Ueberschiebung
\[
\big((k,\Delta)^2,p\big)^2 = \left(\frac{B}{6}\,k,\,p\right)^2.
\]
Denn man hat:
\[
\big((k,\Delta)^2,p\big)^2 = \frac{(k\Delta)^2}{6}\left\{ (kp)^2\Delta_x^2 + (\Delta p)^2 k_x^2 + 4(kp)(\Delta p)\,k_x\Delta_x \right\},
\]
oder, da die beiden ersten Glieder rechts die nämliche Covariante darstellen, wie man unmittelbar erkennt, wenn man $\Delta$ durch $(k_1 k_2)^2 k_{1x}^2 k_{2x}^2$ ersetzt, und da der dritte Term der Klammer nach dem Productsatz sich durch $2(\Delta p)^2 k_x^2 + 2(kp)^2\Delta_x^2 - 2(k\Delta)^2 p$ ersetzen lässt, so wird
\[
\big((k,\Delta)^2,p\big)^2 = \frac{(k\Delta)^2}{6}\left\{ 6(kp)^2\Delta_x^2 - 2(k\Delta)^2 p \right\},
\]
oder
\[
\frac{B}{6}(k,p)^2 = (k\Delta)^2(kp)^2\Delta_x^2 - \frac{1}{3}\,C\,p,
\]
und folglich
\[
2(kp)^2(k\Delta)^2\Delta_x^2 = \frac{1}{3}\,B(k,p)^2 + \frac{2}{3}\,Cp.
\]
Durch Substitution dieses Werthes in (5) erhält man:
\weq{6}{
n - \frac{1}{3}\,Bl = \frac{1}{3}\,B(k,p)^2 + \frac{1}{3}(\Delta,p)^2(A-2P) + \frac{2}{3}\,Cp.
}
Hierin ist $(k,p)^2$ wegen Gleichung (V) bekannt. $(\Delta,p)^2$ aber geht aus der nämlichen Relation (V) durch zweimalige Ueberschiebung über $k$ hervor; dies liefert nämlich:
%CONT>>


% ===== pdf idx 314, printed 300.0 =====
% !PRINTED 300
\weberpage{300}
%<<CONT
\[
(k,l)^2 = m = 2\big(k,(k,p)^2\big)^2 + \frac{1}{3}(k,p)^2(A-2P),
\]
oder:
\[
m = 2(kk_1)^2(k_1p)^2 k_x^2 + \frac{1}{3}(k,p)^2(A-2P).
\]
Der erste Term rechts ist ein Glied der Ueberschiebung $(\Delta,p)^2$; denn man hat:
\[
\begin{aligned}
(\Delta,p)^2 &= (kk_1)^2(k_x^2 k_{1x}^2,\,p)^2 = \frac{(kk_1)^2}{3}\left\{ (k_1p)^2 k_x^2 + 2(kp)(k_1p)\,k_{1x} k_x \right\} \\
&= (kk_1)^2(k_1p)^2 k_x^2 - \frac{1}{3}\,Bp,
\end{aligned}
\]
oder:
\[
2(kk_1)^2(k_1p)^2 k_x^2 = 2(\Delta,p)^2 + \frac{2B}{3}\,p.
\]
Demnach wird:
\weq{7}{
m = 2(\Delta,p)^2 + \frac{2}{3}\,Bp + \frac{A-2P}{3}(k,p)^2.
}
Wir haben also, wenn wir der Einfachheit halber
\[
A - 2P = 6\varrho
\]
setzen, folgende drei Relationen:
\[
\begin{aligned}
\frac{l}{2} &= (k,p)^2 + p\cdot\varrho \quad\text{[vgl.\ (V)]} \\
\frac{m}{2} &= (k,p)^2\cdot\varrho + (\Delta,p)^2 + p\cdot\frac{B}{3} \quad\text{[vgl.\ (7)]} \\
\left(\frac{n}{2}-\frac{B}{6}l\right) &= (k,p)^2\cdot\frac{B}{6} + (\Delta,p)^2\cdot\varrho + p\cdot\frac{C}{3} \quad\text{[vgl.\ (6)].}
\end{aligned}
\]
276. Multiplicirt man die erste Gleichung mit $\varrho^2-\frac{B}{6}$, die zweite mit $-\varrho$, die dritte mit $1$ und addirt, so kommt:
\[
\frac{l}{2}\left(\varrho^2-\frac{B}{6}\right) - \frac{m}{2}\,\varrho + \left(\frac{n}{2}-\frac{B}{6}l\right) = p\left\{ \varrho^3 - \frac{B}{6}\varrho - \frac{B}{3}\varrho + \frac{C}{3} \right\},
\]
oder:
\[
\left(\frac{\varrho^2}{2}-\frac{B}{4}\right) l - \frac{\varrho}{2}\,m + \frac{n}{2} = p\left\{ \varrho^3 - \frac{B}{2}\varrho + \frac{C}{3} \right\}.
\]
Ueberschiebt man endlich diese Relation zweimal über $l$ und ersetzt $A_{ll}$, $A_{ml}$, $A_{nl}$ durch ihre Werthe $2C+\frac{AB}{3}$, resp. $\frac{2}{3}(B^2+AC)$ und $D$, so kommt als zweite Relation für $(lp)^2$:
\[
\left(\frac{\varrho^2}{2}-\frac{B}{4}\right)\left(2C+\frac{AB}{3}\right) - \frac{\varrho}{3}(B^2+AC) + \frac{D}{2} = (lp)^2\left\{ \varrho^3 - \frac{B}{2}\varrho + \frac{C}{3} \right\}.
\]
Ersetzen wir hierin gemäss Gleichung (VI) $(lp)^2$ durch
\[
\frac{2}{3}\,B + \frac{P}{9}(A-2P) = -2\varrho^2 + \frac{A}{3}\varrho + \frac{2}{3}B,
\]
so kommt nach gehöriger Vereinfachung:
%CONT>>


% ===== pdf idx 315, printed 301.0 =====
